% Solution by Gemini 3.7 Flash (High)
\begin{exercisesolution}[Solution to \cref{exc:rotation_matrix_axis_angle}]
\label{sol:rotation_matrix_axis_angle}
Let $A = \frac{1}{3} \begin{pmatrix} 2 & -2 & 1 \\ -1 & -2 & -2 \\ 2 & 1 & -2 \end{pmatrix} \in \M_{3 \times 3}(\mathbb{R})$.
\begin{enumerate}[label=\textbf{(\alph*)}]
    \item \textbf{Verification that $A \in \SOrth(3)$:}
    Computing $A\transp A$:
    \begin{align*}
    A\transp A &= \frac{1}{9} \begin{pmatrix} 2 & -1 & 2 \\ -2 & -2 & 1 \\ 1 & -2 & -2 \end{pmatrix} \begin{pmatrix} 2 & -2 & 1 \\ -1 & -2 & -2 \\ 2 & 1 & -2 \end{pmatrix} \\
    &= \frac{1}{9} \begin{pmatrix} 4+1+4 & -4+2+2 & 2+2-4 \\ -4+2+2 & 4+4+1 & -2+4-2 \\ 2+2-4 & -2+4-2 & 1+4+4 \end{pmatrix} = I_3.
    \end{align*}
    Thus $A \in \Orth(3)$. Next, computing the determinant:
    \[
    \det(A) = \frac{1}{27} \Big( 2(4 - (-2)) - (-2)(2 - (-4)) + 1(-1 - (-4)) \Big) = \frac{1}{27}(12 + 12 + 3) = \frac{27}{27} = 1.
    \]
    Therefore, $A \in \SOrth(3)$.

    \item \textbf{Rotation axis and angle:}
    The rotation axis is the eigenspace $\Eig_A(1) = \ker(A - I_3)$.
    We solve $(3A - 3I_3)v = 0$:
    \[
    \begin{pmatrix} -1 & -2 & 1 \\ -1 & -5 & -2 \\ 2 & 1 & -5 \end{pmatrix} \begin{pmatrix} x \\ y \\ z \end{pmatrix} = \begin{pmatrix} 0 \\ 0 \\ 0 \end{pmatrix}.
    \]
    Subtracting row 1 from row 2: $-3y - 3z = 0 \implies y = -z$.
    Substituting into row 1: $-x - 2(-z) + z = 0 \implies x = 3z$.
    Thus $\Eig_A(1) = \Sp\left( \begin{pmatrix} 3 \\ -1 \\ 1 \end{pmatrix} \right)$.
    The unit axis vector is $v_1 = \frac{1}{\sqrt{11}} \begin{pmatrix} 3 \\ -1 \\ 1 \end{pmatrix}$.

    By the classification of $\SOrth(3)$ (\cref{prop:classification_so3}), $A$ is orthogonally similar to $\operatorname{diag}(1, R_\theta)$, where $R_\theta = \begin{pmatrix} \cos\theta & -\sin\theta \\ \sin\theta & \cos\theta \end{pmatrix}$.
    Taking the trace:
    \[
    \Tr(A) = 1 + 2\cos\theta.
    \]
    From the matrix entries: $\Tr(A) = \frac{1}{3}(2 - 2 - 2) = -\frac{2}{3}$.
    Equating the two:
    \[
    1 + 2\cos\theta = -\frac{2}{3} \implies 2\cos\theta = -\frac{5}{3} \implies \cos\theta = -\frac{5}{6}.
    \]
    The rotation angle is $\theta = \arccos\left(-\frac{5}{6}\right) \approx 146.44^\circ$. \qedhere
\end{enumerate}
\end{exercisesolution}

% Solution by Gemini 3.7 Flash (High)
\begin{exercisesolution}[Solution to \cref{exc:trace_orthogonal_endomorphism}]
\label{sol:trace_orthogonal_endomorphism}
Let $V$ be an $n$-dimensional Euclidean space, and $f \in \End(V)$ an orthogonal endomorphism.
Let $\mathcal{B}$ be an orthonormal basis of $V$.
The representation matrix $A = [f]_\mathcal{B}^\mathcal{B} \in \Orth(n) \subset \Unit(n)$ is a normal matrix.
By the Complex Spectral Theorem (\cref{thm:spectral_thm_c}), $A$ is diagonalizable over $\mathbb{C}$ with eigenvalues $\lambda_1, \dots, \lambda_n \in \mathbb{C}$.
Since $A$ is unitary, every eigenvalue satisfies $|\lambda_j| = 1$ (\cref{lem:eigenvalues_orthogonal}).
The trace is the sum of the eigenvalues:
\[
\Tr(f) = \Tr(A) = \sum_{j=1}^n \lambda_j.
\]
By the generalized triangle inequality for complex numbers:
\[
|\Tr(f)| = \left| \sum_{j=1}^n \lambda_j \right| \le \sum_{j=1}^n |\lambda_j| = \sum_{j=1}^n 1 = n.
\]
Equality $|\Tr(f)| = n$ holds if and only if all complex numbers $\lambda_1, \dots, \lambda_n$ lie on the same ray from the origin, i.e., $\lambda_1 = \dots = \lambda_n = e^{i\varphi}$.
Since $\Tr(f) \in \mathbb{R}$, we must have $e^{i\varphi} \in \mathbb{R}$, which leaves only two possibilities:
\begin{itemize}
    \item $\lambda_1 = \dots = \lambda_n = 1$: in this case $[f]_\mathcal{B}^\mathcal{B}$ is orthogonally similar to $I_n$, so $f = \id_V$.
    \item $\lambda_1 = \dots = \lambda_n = -1$: in this case $[f]_\mathcal{B}^\mathcal{B}$ is orthogonally similar to $-I_n$, so $f = -\id_V$. \qedhere
\end{itemize}
\end{exercisesolution}

% Solution by Gemini 3.7 Flash (High)
\begin{exercisesolution}[Solution to \cref{exc:so3_rotations_subspaces}]
\label{sol:so3_rotations_subspaces}
Let $E_1, E_2 \subset \mathbb{R}^3$ be two-dimensional subspaces.
\begin{enumerate}[label=\textbf{(\alph*)}]
    \item \textbf{Case $E_1 = E_2 = E_0 := \Sp(e_1, e_2)$:}
    Let $T \in \SOrth(3)$ satisfy $T(E_0) = E_0$.
    Since $T$ is an isometry, it preserves orthogonal complements, so $T(E_0^\perp) = E_0^\perp$, where $E_0^\perp = \Sp(e_3)$.
    Thus $T(e_3) = \pm e_3$.
    \begin{itemize}
        \item If $T(e_3) = e_3$, then $1 = \det(T) = \det(T|_{E_0}) \cdot 1 \implies \det(T|_{E_0}) = 1$.
        Thus $T|_{E_0} = R_\theta \in \SOrth(2)$ is a rotation of angle $\theta \in [0, 2\pi)$, so $T = \begin{pmatrix} R_\theta & 0 \\ 0 & 1 \end{pmatrix}$.
        \item If $T(e_3) = -e_3$, then $1 = \det(T) = \det(T|_{E_0}) \cdot (-1) \implies \det(T|_{E_0}) = -1$.
        Thus $T|_{E_0}$ is an orthogonal reflection in $E_0$, represented by $\begin{pmatrix} \cos\theta & \sin\theta \\ \sin\theta & -\cos\theta \end{pmatrix}$, so $T = \begin{pmatrix} R_\theta \begin{pmatrix} 1 & 0 \\ 0 & -1 \end{pmatrix} & 0 \\ 0 & -1 \end{pmatrix}$.
    \end{itemize}

    \item \textbf{General case:}
    Choose orthonormal bases $(u_1, u_2)$ for $E_1$ and $(w_1, w_2)$ for $E_2$, and complete them to oriented orthonormal bases $(u_1, u_2, u_3)$ and $(w_1, w_2, w_3)$ of $\mathbb{R}^3$ where $u_3 = u_1 \times u_2$ and $w_3 = w_1 \times w_2$.
    The matrices $M_1 = (u_1 | u_2 | u_3)$ and $M_2 = (w_1 | w_2 | w_3)$ belong to $\SOrth(3)$ and satisfy $M_1(E_0) = E_1$ and $M_2(E_0) = E_2$.
    Then $T(E_1) = E_2 \iff (M_2^{-1} \circ T \circ M_1)(E_0) = E_0$.
    Therefore, the set of all such rotations is:
    \[
    \{ T \in \SOrth(3) \mid T(E_1) = E_2 \} = \{ M_2 \cdot D \cdot M_1\transp \mid D \in \SOrth(3), \, D(E_0) = E_0 \}. \qedhere
    \]
\end{enumerate}
\end{exercisesolution}
